% ============================================================
%  II. RELATED WORK
% ============================================================
\section{Related Work}
\label{sec:related_work}

\subsection{Robotic Ultrasound Systems}

Robot-assisted ultrasound has been an active research area for over two
decades, driven by the need to reduce operator variability, enable remote
diagnostics, and automate repetitive scanning procedures. Early systems
focused on teleoperation, using robotic arms to replicate the hand movements
of a remote sonographer~\cite{li2021overview}. More recent work has moved
toward increasing levels of autonomy, with robots capable of autonomously
positioning and maintaining contact of the probe on the patient's body.

A fundamental challenge in autonomous ultrasound is probe--skin contact
control. Gilbertson and Anthony~\cite{gilbertson2015} proposed a force and
position control system for freehand ultrasound that decouples normal contact
force from lateral probe positioning, establishing a framework widely adopted
by subsequent systems. Building on this, Dyck et al.~\cite{dyck2022} extended
impedance control to arbitrary curved surfaces using discrete differential
geometry, enabling consistent probe orientation on non-planar anatomy---a
capability directly relevant to abdominal scanning.

Autonomous scan planning has been tackled from multiple angles. Suligoj
et al.~\cite{suligoj2021} presented RobUSt, a system that autonomously
acquires medical images using a pre-defined anatomical protocol, while Duan
et al.~\cite{duan2022} demonstrated optimisation-based control for scoliosis
assessment with variable impedance. A comprehensive overview of techniques for
autonomous robotic ultrasound acquisitions is provided by Li et
al.~\cite{li2021overview}, and a broader review of robot-assisted ultrasound
systems is given in~\cite{du2024review}.

All of the above systems rely on fixed-base manipulators and require the
patient to be positioned within the arm's reachable workspace. In our own
prior work~\cite{farsoni2025autonomous}, we demonstrated that a fixed
collaborative robot arm can autonomously screen the abdominal aorta using
impedance control and deep learning segmentation. The present paper extends
this line of work to a fully mobile legged platform, removing the
pre-positioning requirement and enabling operation in unstructured environments.

\subsection{Whole-Body Control for Legged Manipulators}

\ac{WBC} has emerged as the dominant paradigm for coordinating locomotion and
manipulation on legged robots. Khatib~\cite{khatib1987} introduced the
operational space formulation that underpins most modern \ac{WBC} approaches,
enabling task-space control through the robot's full kinematic chain.
Yoshikawa's manipulability measure~\cite{yoshikawa1985} provides a scalar
index of arm dexterity that has been widely used to avoid singular
configurations.

For mobile manipulators, Bayle et al.~\cite{bayle2003} analysed the
manipulability of wheeled platforms coupled with robot arms, showing that base
motion can be used to maintain the arm in high-dexterity configurations.
Haviland et al.~\cite{haviland2022} extended this to a holistic reactive
formulation that simultaneously optimises base and arm velocities through a
single \ac{QP}, achieving smooth coordinated motion on wheeled platforms.

On legged robots, Bellicoso et al.~\cite{bellicoso2019} presented ALMA,
demonstrating articulated locomotion and manipulation with a torque-controlled
legged platform. Minniti et al.~\cite{minniti2019} proposed a whole-body
\ac{MPC} framework for dynamically stable mobile manipulation, and Sleiman
et al.~\cite{sleiman2021} unified locomotion and manipulation into a single
\ac{MPC} problem. These approaches have been validated primarily on ANYmal
with structured targets at known locations. The Boston Dynamics Spot platform
has been used for manipulation tasks with an external arm by Zimmermann
et al.~\cite{zimmermann2021}, who demonstrated dynamic grasping. However,
none of the existing legged \ac{WBC} systems have addressed
perception-driven, multi-phase medical scanning tasks where the target must
be detected and tracked throughout the mission.

\subsection{Visual Perception for Patient Localisation}

Accurate patient localisation from \ac{RGBD} data is a prerequisite for
autonomous mobile scanning. Skeleton-based approaches, which estimate 3D joint
positions from monocular or depth-augmented images, have matured significantly
with deep learning. YOLO-Pose~\cite{maji2022} extended the \ac{YOLO}
detection framework to multi-person pose estimation using an object keypoint
similarity loss, achieving real-time performance suitable for robot perception
pipelines.

In our system, \ac{YOLO}-based skeleton detection is combined with a \ac{KF}
for 3D torso tracking, providing robust localisation even under partial
occlusion and during platform motion. The \ac{KF} predicts torso position
between detections and smooths noisy depth measurements from the Orbbec
\ac{RGBD} camera, enabling stable approach trajectories without requiring the
robot to stop for reliable detection.

In summary, while individual components---robotic ultrasound, \ac{WBC}, and
visual patient tracking---have each received significant attention, no prior
work has integrated them into a unified system capable of autonomously
performing a medical scan on a legged platform navigating to an unknown
patient location. The present paper addresses this gap.
